\chapter{Chakshu: the same machine for one thermal band}\label{ch:chakshu}

\begin{onething}
Chakshu, the product's thermal model, is the SegFormer masked autoencoder of the last nine chapters with one
input band of surface temperature in degrees Celsius, widths shrunk to 16, 32, 64 and 96, one block per stage and a
feed-forward expansion of 2: 406,496 trainable numbers that you can count by hand, trained to fill in 70\% hidden
tokens of 256-pixel Landsat patches, and scored by the absolute temperature error; its one serious weakness is
outside the network, in the normalisation path that was meant to let it read HotSat-1 and does not.
\end{onething}

You have now built Indradhanu from its parts. This chapter uses every one of those parts again, on a simpler input,
and asks you to redo the arithmetic at new sizes. Nothing here is new machinery; what is new is the thermal domain,
the training recipe and three places where the code and its documentation part company. It is also the chapter
where you can check yourself: if a shape or a count below surprises you, the chapter that explains it is named
beside it.

Provenance tags abbreviate paths. \file{components/} is \file{app/foundation_models/components/}; the shipped
checkpoint is \file{segformer_mae_v0.8.0_epoch495.pt} in the stack's \file{/models/segformer_mae/} volume, which is
not in the repository. Its configuration is also in the repository as \file{configs/seg_former_exp_8.json}, which
matches the checkpoint's saved configuration key for key. Every count and shape tagged \tagC{} comes from
\file{research/textbook/tools/compute_chakshu.py}, which loads that checkpoint (copied from the volume on 2026-10-31)
into the real classes; its output is saved beside it as \file{compute_chakshu_output.txt}.

\begin{analogy}
A town has a thermometer on every street corner. One morning you cover seven out of every ten and ask an
apprentice to say what each covered one reads, from the uncovered ones around it. On an ordinary morning the
apprentice is right to a fraction of a degree, because temperature varies smoothly across a town. If a warehouse is
burning, the covered thermometer beside it reads 20 degrees more than anything its neighbours suggest, and the
apprentice's guess is 20 degrees wrong. That error is the alarm. Two details of the analogy will matter later: the
apprentice must never see the covered thermometer itself, and if the fire is large enough to reach an uncovered
thermometer too, the apprentice sees part of the fire and the alarm is quieter.
\end{analogy}

%======================================================================
\section{Why a thermal model, and why this one}\label{sec:chakshu:why}

Two of the product's five sensors are \emph{thermal} (\cref{ch:light,ch:sensors}): Landsat~9's band~10 and
SatVu's HotSat-1. Each records one number per pixel, related to how much heat the ground radiates, near
$10.9\,\mu$m (the product files thermal scenes under the wavelength \code{10900.0} nm,
\file{backend/allotrope/action_types/_anomaly_scoring_run.py}). A thermal anomaly is something hotter or colder than
its surroundings say it should be: a fire, a furnace, a cooling-water plume, a vessel's engine.

\begin{definition}[Land surface temperature]
The \emph{land surface temperature}\index{land surface temperature} of a pixel is the temperature of the ground's
skin, as inferred from the radiation it emits. The product stores it in degrees Celsius for Landsat~9.
\end{definition}

The product has two ways to score a thermal scene. The classical \code{thermal_grx} detector (\cref{ch:rx}) measures
how far each pixel's temperature sits from the scene's mean, in standard deviations: a hot pixel in a cold field
scores high, but so does every ordinary hot roof in a city. A reconstruction model asks a sharper question, ``is
this pixel hot \emph{for where it is}?'', because its prediction comes from the pixel's own neighbourhood
(\cref{ch:recon}).

Which reconstruction model? \Cref{ch:convae} ran the five convolutional autoencoders over a smooth synthetic field with the
product's own inference: their median background error ran from 0.58\,°C (Drashta) to 10.9\,°C (Pratibimba, which
never trained with hidden pixels). Three of them cannot tell a hidden pixel from a real one at 0\,°C. The SegFormer
design avoids that: hidden tokens are \emph{deleted} before the encoder (\cref{ch:masking}), so there is no stand-in
value to misread, and the decoder's pixel shuffle gives every pixel its own output (\cref{ch:decoder}), so a
one-pixel hot spot is not averaged into its neighbours. On the same synthetic field Chakshu's background error is 0.019\,°C (\cref{sec:chakshu:score}).

\paragraph{Why so small.} Indradhanu's encoder is 32, 64, 160, 256 channels wide, with two blocks per stage.
Chakshu's is 16, 32, 64, 96 with one. A thermal pixel carries one number, not 165, so there is far less to model;
and the project's own record says the small model won. The manifest lists the larger SegFormer-B0 run (v0.1.0,
3,714,086 numbers) as an alternative with validation loss 6.59877, against 0.2565 for the small one, and notes that
the small one ``currently outperforms it 25$\times$ on val loss'' \tagMr{the \texttt{/models/segformer\_mae/current.json}
manifest}. Read that comparison with care: the large run had trained for 10 epochs, the small one for 495, and the
two losses were measured on different patch sizes (128 and 256 pixels). It shows that the large run was unfinished,
not that the large architecture is worse.

%======================================================================
\section{The input: one band, in Celsius}\label{sec:chakshu:input}

\subsection{From Landsat's numbers to degrees}

Landsat~9's Level-2 surface-temperature product stores band~10 as a 16-bit integer, a \emph{digital
number}\index{digital number} (DN). The builder (\file{app/utils/dataset_builder/landsat_dataset_builder.py}) converts
it with the published scale and offset (\file{app/utils/data_transformations/l2sp_dn_to_temperature_transformer.py}):
\begin{equation}\label{eq:chakshu:dn}
T_{\mathrm{K}} = 0.00341802\cdot\mathrm{DN} + 149.0, \qquad T = T_{\mathrm{K}} - 273.15 .
\end{equation}
The unit is fixed to Celsius in the builder; it is not an option.

\begin{worked}[label={we:chakshu:dn}]{Two Landsat pixels into degrees}
An ordinary pixel, DN $=43{,}500$:
\begin{align*}
0.00341802\times43{,}500 &= 148.68387 && \text{scale}\\
148.68387+149.0 &= 297.68387\ \text{K} && \text{offset}\\
297.68387-273.15 &= 24.53\ \text{°C} && \text{Kelvin to Celsius}
\end{align*}
A hot pixel, DN $=52{,}000$: $0.00341802\times52{,}000=177.73704$; $+149.0=326.73704$\,K; $-273.15=53.59$\,°C.
\end{worked}

\subsection{The normalisation the checkpoint carries}

Chakshu normalises its input with one mean and one spread, stored in the checkpoint as buffers
(\cref{ch:bandpipeline}): $\mu_T=24.5756$ and $\sigma_T=13.5744$ \tagMr{checkpoint buffers \texttt{normalize.mean},
\texttt{normalize.std}}. The same two numbers are in \file{app/constants/thermal_pixel_consts.json}, measured over
10,000 Landsat patches of 64 pixels (25,413,600 valid pixels) \tagM{app/constants/thermal_pixel_consts.json}.
\begin{equation}\label{eq:chakshu:z}
z = \frac{T-\mu_T}{\sigma_T} .
\end{equation}

\begin{worked}[label={we:chakshu:z}]{Normalising the two pixels, and an invalid one}
\begin{align*}
\text{ordinary} &: (24.5339-24.5756)/13.5744 = -0.0417/13.5744 = -0.0031\\
\text{hot} &: (53.5870-24.5756)/13.5744 = 29.0114/13.5744 = 2.1372\\
\text{invalid, zeroed to } 0\text{ °C} &: (0-24.5756)/13.5744 = -24.5756/13.5744 = -1.8104
\end{align*}
The model zeroes invalid pixels \emph{before} normalising (\lstinline!x = x * mask! in \file{components/seg_former_mae.py}),
so a cloud pixel arrives as $-1.81$: a consistent value that the network learns to read as ``no measurement''
(\cref{ch:bandpipeline}) \tagC.
\end{worked}

\begin{sourcenote}[: Celsius near 25, not Kelvin near 290]
Three places describe Chakshu's input wrongly. The comment above the per-scene override in
\file{backend/allotrope/action_types/_anomaly_scoring_run.py} says the models assume ``Kelvin / reflectance'' (the same
comment, ten lines later, says Landsat is ``already in Celsius''). \code{PixelStatsOverride}'s docstring
(\file{app/models/training/inference_config.py}) and \file{docs/lld/pixel-stats.md} say a model ``fitted on Celsius
(\textasciitilde290 +/- 10)'', and the LLD and the course (part AD10) work the HotSat example with a mean of 290 and
a spread of 10. 290 is a Kelvin-sized number. The code and the checkpoint agree on Celsius: the vendable's field is
documented as ``surface temperatures in celsius'' (\file{app/models/dataset/vendables.py}), the builder converts to
Celsius, and the buffers are $24.58\pm13.57$. In Kelvin the mean would be $24.5756+273.15=297.73$.
The comment in \file{components/spatial_decoder.py} (``e.g.\ surface temperatures in Kelvin'') is the same slip in
the convolutional family.
\end{sourcenote}

%======================================================================
\section{The configuration}\label{sec:chakshu:config}

\begin{table}[htbp]
\centering\small
\begin{tabular}{lll}
\toprule
setting & Chakshu (v0.8.0, served) & class default (SegFormer-B0)\\
\midrule
\texttt{in\_channels} & 1 & 1\\
\texttt{embed\_dims} & 16, 32, 64, 96 & 32, 64, 160, 256\\
\texttt{num\_heads} & 1, 2, 4, 8 & 1, 2, 5, 8\\
channels per head & 16, 16, 16, 12 & 32, 32, 32, 32\\
\texttt{reduction\_ratios} & 8, 4, 2, 1 & 8, 4, 2, 1\\
\texttt{num\_blocks} & 1, 1, 1, 1 & 2, 2, 2, 2\\
\texttt{expansion\_ratio} $E$ & 2 & 4\\
\texttt{decoder\_dim} & 96 & 256\\
\texttt{drop\_rate} & 0.3 (training) & 0.3\\
\texttt{mask\_ratio} $\rho$ & 0.7 & 0.5\\
\texttt{trim\_fraction} $\tau$ & 0.05 & 0.0\\
patch size in training & 256 & --\\
\bottomrule
\end{tabular}
\caption{Chakshu against the defaults of \texttt{SegFormerMAEConfig} (\texttt{app/models/training/training\_config.py}).
Source: \tagM{configs/seg_former_exp_8.json}, identical to the checkpoint's saved configuration \tagC.}
\label{tab:chakshu:config}
\end{table}

\begin{sourcenote}[: the served model is not SegFormer-B0]
\file{research/model_break_down/04_segformer_mae_thermal.md} draws its architecture diagram with stage widths 32, 64,
160, 256 and two blocks per stage, and says ``The default config (used by both checkpoints)'' is SegFormer-B0. Only
the unused alternative, v0.1.0, has that configuration. The served v0.8.0 is the shrunk one of
\cref{tab:chakshu:config}; the same document's checkpoint table, further down, says so correctly. The book's earlier
remark that Chakshu ``shares almost all its code with Indradhanu'' (\cref{ch:light}) is true of the classes, not of
the sizes.
\end{sourcenote}

How does the product build this model? The worker never opens the checkpoint's configuration. The resolver
rebuilds it from the manifest (\code{_model_config_from_manifest} in
\file{backend/allotrope/foundation_models/resolver.py}): widths from the manifest's \code{encoder_dims} string
\lstinline!"[16, 32, 64, 96]"!, decoder width from \code{decoder_dim}, and \emph{fixed} heads $[1,2,4,8]$, ratios
$[8,4,2,1]$, one block per stage and expansion 2, hard-coded for \code{segformer_mae}. Those happen to equal the
checkpoint's, and the strict load succeeds \tagC. A future thermal checkpoint with two blocks per stage would build
the wrong model and fail to load.

%======================================================================
\section{One forward pass at full size}\label{sec:chakshu:trace}

Follow one masked pass, as in training, with patches of $256\times256$ pixels, all valid, and the training mask
ratio 0.7. Training batches hold 64 patches; take $B=8$, the product's inference batch, since the batch size changes
no shape except the first axis.

\subsection{Stage 1: the patch embedding and the deletion}

The stage-1 patch embedding is a convolution with kernel 4, stride 4 and no padding (\cref{ch:patchembed}) from 1
channel to 16. A $256\times256$ patch becomes a $64\times64$ grid:
\begin{align*}
\left\lfloor\frac{256+2\cdot0-4}{4}\right\rfloor+1 &= \lfloor 63\rfloor+1 = 64 && \text{output-size formula, \cref{ch:shapes}}\\
64\times64 &= 4{,}096 \text{ tokens of 16 numbers.}
\end{align*}
Each token is a weighted sum of its own $4\times4=16$ temperatures and nothing else: the weights form a $16\times16$
matrix, one row per output channel. With one input band, stage~1 is exactly a linear map from a block's 16
temperatures to 16 features, followed by LayerNorm.

\begin{worked}[label={we:chakshu:mask}]{How many tokens a training pass hides}
The hiding count rounds down (\cref{ch:masking}):
\begin{align*}
4{,}096\times0.7 &= 2{,}867.2 && \text{tokens to hide, unrounded}\\
\lfloor2{,}867.2\rfloor &= 2{,}867 && \text{the code's \texttt{.long()} rounds down}\\
4{,}096-2{,}867 &= 1{,}229 && \text{tokens the encoder sees}\\
2{,}867\times16 &= 45{,}872 && \text{hidden pixels, the loss's positions}
\end{align*}
Running the real encoder gives 2,867 hidden and 1,229 kept per item \tagC.
\end{worked}

The encoder sees 1,229 of 4,096 tokens, 30\%. Indradhanu hides 65\% (\cref{ch:masking}); Chakshu hides more because a
temperature field is smoother than a reflectance cube, so the same context predicts more of it. That is a design
judgement recorded only in the configuration's value; no document gives the reason.

\subsection{The four stages}

\begin{table}[htbp]
\centering\small
\begin{tabular}{@{}lp{0.36\linewidth}lp{0.2\linewidth}@{}}
\toprule
stage & step & shape after the step & note\\
\midrule
input & normalise, $z=(T-\mu_T)/\sigma_T$ & \shape{(8, 1, 256, 256)} & \cref{we:chakshu:z}\\
1 & patch embed $k4\,s4\,p0$, $1\to16$ & \shape{(8, 4096, 16)} & 64 $\times$ 64 grid\\
 & remove hidden tokens & \shape{(8, 1229, 16)} & 2,867 deleted\\
 & 1 block: attention, Mix-FFN ($16\to32\to16$) & \shape{(8, 1229, 16)} & full attention; no depthwise conv\\
 & LayerNorm, restore, reshape & \shape{(8, 16, 64, 64)} = F1 & zeros at hidden tokens\\\midrule
2 & patch embed $k3\,s2\,p1$, $16\to32$ & \shape{(8, 1024, 32)} & 32 $\times$ 32 grid\\
 & 1 block, $R=4$, 2 heads; FFN $32\to64\to32$ & \shape{(8, 1024, 32)} & 64 keys\\
 & LayerNorm, reshape & \shape{(8, 32, 32, 32)} = F2 & \\\midrule
3 & patch embed, $32\to64$; 1 block, $R=2$, 4 heads & \shape{(8, 256, 64)} & 64 keys\\
 & reshape & \shape{(8, 64, 16, 16)} = F3 & \\\midrule
4 & patch embed, $64\to96$; 1 block, $R=1$, 8 heads & \shape{(8, 64, 96)} & 64 keys\\
 & reshape & \shape{(8, 96, 8, 8)} = F4 & \\
\bottomrule
\end{tabular}
\caption{Chakshu's encoder on a masked training pass, $B=8$, $256\times256$, $\rho=0.7$. Forward hooks on the real
modules record every shape \tagC.}
\label{tab:chakshu:encoder}
\end{table}

Two rows deserve a second look. At stage~1 the token count after deletion, 1,229, is not $64\times64$, so efficient
self-attention cannot rebuild its grid and attends over all kept tokens with no reduction, and Mix-FFN skips its
depthwise convolution (\cref{ch:esa,ch:block}, the \code{can_reduce} test). And at stage~4 each head is
$96/8=12$ channels wide, not a power of two: harmless, since nothing in attention needs one.

\begin{worked}[label={we:chakshu:keys}]{Every stage attends to 64 keys}
Efficient self-attention shrinks the key grid by $R$ on each side (\cref{ch:esa}). At $256\times256$:
\begin{align*}
\text{stage 1} &: (64/8)^2 = 8^2 = 64 \text{ keys for } 4{,}096 \text{ queries}\\
\text{stage 2} &: (32/4)^2 = 8^2 = 64 \text{ keys for } 1{,}024 \text{ queries}\\
\text{stage 3} &: (16/2)^2 = 8^2 = 64 \text{ keys for } 256 \text{ queries}\\
\text{stage 4} &: (8/1)^2 = 64 \text{ keys for } 64 \text{ queries}
\end{align*}
Unmasked, stage~1's score table has $4{,}096\times64=262{,}144$ entries per head. Masked, it falls back to full
attention among the kept tokens: in training $1{,}229^2=1{,}510{,}441$, and at inference, where half the tokens are
kept, $2{,}048^2=4{,}194{,}304$, sixteen times the unmasked table \tagC. The fallback is the expensive case.
\end{worked}

\subsection{The decoder}

The decoder is \code{SegFormerDecoder} (\cref{ch:decoder}) at width 96, writing $1\times4\times4=16$ channels for the
pixel shuffle:
\begin{center}\small
\begin{tabular}{@{}lll@{}}
\toprule
step & operation & shape\\
\midrule
project & $1\times1$ conv of each F$_i$ to 96 channels & F1 \shape{(8, 96, 64, 64)}, \dots, F4 \shape{(8, 96, 8, 8)}\\
upsample & bilinear, to $64\times64$ & four of \shape{(8, 96, 64, 64)}\\
concatenate & along channels & \shape{(8, 384, 64, 64)}\\
fuse & $1\times1$ conv $384\to96$, GELU & \shape{(8, 96, 64, 64)}\\
refine & $3\times3$ conv $96\to96$, GELU & \shape{(8, 96, 64, 64)}\\
sub-pixel & $3\times3$ conv $96\to16$, no activation & \shape{(8, 16, 64, 64)}\\
pixel shuffle & factor 4 & \shape{(8, 1, 256, 256)}\\
denormalise & $T=z\,\sigma_T+\mu_T$ & \shape{(8, 1, 256, 256)}, in °C\\
\bottomrule
\end{tabular}
\end{center}
Every shape was recorded by forward hooks \tagC. \Cref{fig:chakshu:trace} draws the whole pass.

\begin{figure}[htbp]\centering
\begin{tikzpicture}[node distance=3.6mm,
  box/.style={draw,rounded corners=2pt,minimum width=5.6cm,minimum height=6mm,font=\small,align=center},
  lab/.style={font=\footnotesize\ttfamily,anchor=west},
  side/.style={font=\footnotesize,anchor=east}]
\node[box,fill=loss!12] (in) {temperature $T$, °C};
\node[box,below=of in] (mask) {zero invalid pixels; $z=(T-24.58)/13.57$};
\node[box,below=of mask,fill=spat!15] (s1) {stage 1: embed $1\to16$, delete 70\%, 1 block};
\node[box,below=of s1,fill=spat!15] (s2) {stage 2: embed $16\to32$, 1 block, $R=4$};
\node[box,below=of s2,fill=spat!15] (s3) {stage 3: embed $32\to64$, 1 block, $R=2$};
\node[box,below=of s3,fill=spat!15] (s4) {stage 4: embed $64\to96$, 1 block, $R=1$};
\node[box,below=of s4,fill=attn!12] (fuse) {project to 96, upsample, fuse $384\to96$};
\node[box,below=of fuse,fill=attn!12] (ref) {refine $3\times3$; sub-pixel $96\to16$};
\node[box,below=of ref,fill=attn!12] (ps) {pixel shuffle $\times4$};
\node[box,below=of ps,fill=loss!12] (out) {denormalise: reconstruction $\hat T$, °C};
\foreach \a/\b in {in/mask,mask/s1,s1/s2,s2/s3,s3/s4,s4/fuse,fuse/ref,ref/ps,ps/out}
  \draw[-{Stealth}] (\a) -- (\b);
\node[lab] at ($(in.east)+(2mm,0)$) {(8, 1, 256, 256)};
\node[lab] at ($(s1.east)+(2mm,0)$) {(8, 4096, 16) -> 1229 kept};
\node[lab] at ($(s1.east)+(2mm,-3.4mm)$) {F1 (8, 16, 64, 64)};
\node[lab] at ($(s2.east)+(2mm,0)$) {F2 (8, 32, 32, 32)};
\node[lab] at ($(s3.east)+(2mm,0)$) {F3 (8, 64, 16, 16)};
\node[lab] at ($(s4.east)+(2mm,0)$) {F4 (8, 96, 8, 8)};
\node[lab] at ($(fuse.east)+(2mm,0)$) {(8, 96, 64, 64)};
\node[lab] at ($(ref.east)+(2mm,0)$) {(8, 16, 64, 64)};
\node[lab] at ($(ps.east)+(2mm,0)$) {(8, 1, 256, 256)};
\node[lab] at ($(out.east)+(2mm,0)$) {(8, 1, 256, 256)};
\node[side] at ($(mask.west)+(-2mm,0)$) {ch.~\ref{ch:bandpipeline}};
\node[side] at ($(s1.west)+(-2mm,0)$) {ch.~\ref{ch:patchembed}, \ref{ch:masking}};
\node[side] at ($(s2.west)+(-2mm,0)$) {ch.~\ref{ch:esa}--\ref{ch:encoder}};
\node[side] at ($(fuse.west)+(-2mm,0)$) {ch.~\ref{ch:decoder}};
\end{tikzpicture}
\caption{Chakshu's forward pass on a masked batch of eight $256\times256$ Landsat patches, with the shape after each
box (recorded by forward hooks on the real model \tagC). Blue: encoder; plum: decoder; amber: the two ends, in
degrees Celsius. The chapters on the left explain each box.}
\label{fig:chakshu:trace}
\end{figure}

%======================================================================
\section{The parameters, counted by hand}\label{sec:chakshu:params}

\Cref{ch:params} counted a block with expansion 4. Chakshu's expansion is 2, so first generalise the block formula to
any expansion $E$, then use it. The counting rules are those of \cref{sec:params:rules}: a linear layer from $a$ to
$b$ has $ab+b$, a convolution $c_\text{in}c_\text{out}k^2+c_\text{out}$, a depthwise convolution $ck^2+c$, a
LayerNorm $2c$.

\subsection{A block of width \texorpdfstring{$C$}{C}, reduction \texorpdfstring{$R$}{R}, expansion \texorpdfstring{$E$}{E}}

\begin{align*}
\text{norm1} &: 2C\\
\text{q, k, v, proj} &: 4\,(C\cdot C+C) = 4C^2+4C\\
\text{reduction conv } (R>1) &: C\cdot C\cdot R^2+C = R^2C^2+C\\
\text{reduction norm } (R>1) &: 2C\\
\text{norm2} &: 2C\\
\text{fc1 } (C\to EC) &: C\cdot EC+EC = EC^2+EC\\
\text{depthwise } 3\times3 \text{ on } EC &: EC\cdot9+EC = 10EC\\
\text{fc2 } (EC\to C) &: EC\cdot C+C = EC^2+C
\end{align*}
Collect the $C^2$ terms: $4+R^2+E+E=4+R^2+2E$. Collect the $C$ terms: $2+4+1+2+2+E+10E+1=12+11E$. So
\begin{equation}\label{eq:chakshu:block}
P_\text{block}(C,R,E) = (4+2E+R^2)\,C^2 + (12+11E)\,C \quad (R>1),
\end{equation}
and for $R=1$ drop the reduction convolution and its norm: $(4+2E)\,C^2+(9+11E)\,C$.

\emph{Check against \cref{ch:params}.} With $E=4$ the coefficients are $4+8+R^2=12+R^2$ and $12+44=56$.
\Cref{eq:params:block} gives $12C^2+53C+R^2C^2+3C=(12+R^2)C^2+56C$ for $R>1$: the two agree. With $E=2$ the
coefficients are $8+R^2$ and $34$ (31 for $R=1$).

\begin{worked}[label={we:chakshu:blocks}]{Chakshu's four blocks}
\begin{align*}
C=16,R=8 &: (8+64)\times256+34\times16 = 72\times256+544 = 18{,}432+544 = 18{,}976\\
C=32,R=4 &: (8+16)\times1{,}024+34\times32 = 24{,}576+1{,}088 = 25{,}664\\
C=64,R=2 &: (8+4)\times4{,}096+34\times64 = 49{,}152+2{,}176 = 51{,}328\\
C=96,R=1 &: 8\times9{,}216+31\times96 = 73{,}728+2{,}976 = 76{,}704
\end{align*}
Line by line for stage 1, to see where the weight is: norm1 32; q, k, v, proj $4\times272=1{,}088$; reduction
$16\times16\times64+16=16{,}400$; reduction norm 32; norm2 32; fc1 $16\times32+32=544$; depthwise $32\times9+32=320$;
fc2 $32\times16+16=528$. Sum: $32+1{,}088+16{,}400+32+32+544+320+528=18{,}976$. The $8\times8$ reduction convolution
is $16{,}400/18{,}976=86\%$ of the block. All four match the built model \tagC.
\end{worked}

\subsection{The rest of the encoder}

Patch embeddings, each a convolution plus a LayerNorm of $2c_\text{out}$:
\begin{align*}
\text{stage 1 } (k=4) &: 1\times16\times16+16+32 = 256+48 = 304\\
\text{stage 2 } (k=3) &: 16\times32\times9+32+64 = 4{,}608+96 = 4{,}704\\
\text{stage 3} &: 32\times64\times9+64+128 = 18{,}432+192 = 18{,}624\\
\text{stage 4} &: 64\times96\times9+96+192 = 55{,}296+288 = 55{,}584\\
\text{sum} &: 304+4{,}704 = 5{,}008;\ +18{,}624 = 23{,}632;\ +55{,}584 = 79{,}216
\end{align*}
Stage-end LayerNorms: $2\times(16+32+64+96)=2\times208=416$. Blocks: $18{,}976+25{,}664=44{,}640$;
$+51{,}328=95{,}968$; $+76{,}704=172{,}672$. Encoder:
\[
172{,}672+79{,}216+416 = 252{,}304 .
\]

\subsection{The decoder and the total}

\begin{align*}
\text{projections} &: (16+32+64+96)\times96+4\times96 = 208\times96+384 = 19{,}968+384 = 20{,}352\\
\text{fuse} &: 384\times96+96 = 36{,}864+96 = 36{,}960\\
\text{refine } 3\times3 &: 96\times96\times9+96 = 82{,}944+96 = 83{,}040\\
\text{sub-pixel } 3\times3 &: 96\times16\times9+16 = 13{,}824+16 = 13{,}840\\
\text{decoder} &: 20{,}352+36{,}960 = 57{,}312;\ +83{,}040 = 140{,}352;\ +13{,}840 = 154{,}192
\end{align*}
\begin{equation}\label{eq:chakshu:total}
252{,}304+154{,}192 = 406{,}496 \text{ trainable numbers.}
\end{equation}
\lstinline!sum(p.numel() for p in model.parameters())! on the loaded checkpoint returns 406,496 \tagC.

\begin{sourcenote}[: 406,500 counts four buffers]
The manifest, \file{docs/04-models.md} (``406 k params'') and the course (AD10: ``406,500 parameters'') quote 406,500.
That is the number of values in the checkpoint's \code{state_dict}, which also holds four buffers: the normalisation
mean and spread, once in \code{PixelNormalize} and once in \code{PixelDenormalize}. $406{,}496+4=406{,}500$. The
same distinction explained 5,506,629 against 5,507,354 for Indradhanu (\cref{ch:recon}).
\end{sourcenote}

\begin{figure}[htbp]\centering
\begin{tikzpicture}
\begin{axis}[xbar, bar shift=0pt, width=0.8\linewidth, height=6.2cm, xmin=0, xmax=100000,
  symbolic y coords={sub-pixel,refine,fuse,projections,patch embeds,stage 4 block,stage 3 block,stage 2 block,stage 1 block},
  ytick={sub-pixel,refine,fuse,projections,patch embeds,stage 4 block,stage 3 block,stage 2 block,stage 1 block},
  yticklabel style={font=\small}, xlabel={trainable parameters},
  scaled x ticks=false, xticklabel style={/pgf/number format/fixed,/pgf/number format/1000 sep={{,}},font=\scriptsize},
  nodes near coords, nodes near coords style={font=\scriptsize,/pgf/number format/1000 sep={{,}}},
  bar width=3.4mm, enlarge y limits=0.07]
\addplot[fill=spat!40,draw=spat] coordinates {(18976,stage 1 block) (25664,stage 2 block) (51328,stage 3 block)
  (76704,stage 4 block) (79216,patch embeds)};
\addplot[fill=attn!35,draw=attn] coordinates {(20352,projections) (36960,fuse) (83040,refine) (13840,sub-pixel)};
\end{axis}
\end{tikzpicture}
\caption{Where Chakshu's 406,496 parameters live, drawn to scale from \cref{we:chakshu:blocks} and the sums above
(blue encoder, plum decoder). The largest single layer is the decoder's refine convolution, 20\% of the model; the
stage-4 patch embedding (55,584) alone outweighs the first three blocks together (stage-end norms, 416, not drawn).}
\label{fig:chakshu:params}
\end{figure}

\subsection{The unused alternative, and why it cannot load}

The manifest's alternative, v0.1.0, is SegFormer-B0 with 3,714,086 values in its \code{state_dict}. Today's code
cannot load it. Its stage-1 kernel is $7\times7$ (\lstinline!(32, 1, 7, 7)!), from before the change to non-overlapping
$4\times4$ tokens (\cref{ch:patchembed}), and its decoder has a BatchNorm after the fuse and a single final
projection where today's has the two refine convolutions. A strict load reports four missing keys, seven unexpected
ones and one size mismatch \tagC. Selecting it in the product would fail at the load.

%======================================================================
\section{Training: 70\% hidden, the trimmed L1}\label{sec:chakshu:train}

The trainer is \file{app/foundation_models/trainers/segformer_mae_trainer.py}. Each step:
\begin{enumerate}
\item builds the validity mask as \lstinline!pure_validity * predicted_cloud_mask!: the provider's valid pixels and the
  cloud mask of the scene's own temperature mixture model (\cref{sec:product:thermal});
\item drops patches with less than 40\% valid pixels, the rule Indradhanu's trainers share (\cref{sec:losses:forty}
  covers it later);
\item turns pixel validity into token validity (a token is valid if more than half of its 16 pixels are) and hides
  $\lfloor0.7\times\text{valid tokens}\rfloor$ of the valid ones at random (\cref{ch:masking});
\item runs the model with that keep mask and takes the absolute error at every hidden, valid pixel;
\item trims, then averages.
\end{enumerate}

\subsection{The trimmed L1}

\begin{definition}[Trimmed mean]
The \emph{trimmed mean}\index{trimmed loss} of $n$ numbers at trim fraction $\tau$ sorts them, keeps the smallest
$\lfloor n(1-\tau)\rfloor$ and averages those. Chakshu's loss is the trimmed mean of the per-pixel absolute errors
$|\hat T-T|$ over the hidden, valid pixels of the whole batch, with $\tau=0.05$.
\end{definition}

The code says why: ``The highest-error pixels are assumed to be anomalies - dropping them from the training signal
teaches the model to reconstruct only normal background patterns.'' The training archive has no anomaly labels, so
any fire in it would otherwise be a target the model is graded on reconstructing, which is exactly what you do not
want it to learn (\cref{sec:recon:l1l2}).

\begin{worked}[label={we:chakshu:trim}]{Trimming 20 errors at 5\%}
Twenty hidden pixels have errors (°C) of 0.10, 0.20, 0.15, 0.30, 0.05, 0.25, 0.20, 0.10, 0.15, 0.10, 0.20, 0.30,
0.25, 0.10, 0.15, 0.20, 0.05, 0.10, 0.20, and one fire pixel at 9.00 \tagA.
\begin{align*}
\text{keep} &: \lfloor20\times(1-0.05)\rfloor = \lfloor19.0\rfloor = 19 \text{ smallest}\\
\text{sum of all 20} &: 3.15+9.00 = 12.15 \quad(\text{the nineteen small ones sum to }3.15)\\
\text{plain mean} &: 12.15/20 = 0.6075\\
\text{trimmed mean} &: 3.15/19 = 0.1658
\end{align*}
To get 3.15: $2\times0.05=0.10$; $5\times0.10=0.50$; $3\times0.15=0.45$; $5\times0.20=1.00$; $2\times0.25=0.50$;
$2\times0.30=0.60$; and $0.10+0.50+0.45+1.00+0.50+0.60=3.15$. One fire pixel had made the loss almost four times
larger. After the trim it contributes nothing, and its gradient is zero: the model is not pushed to reconstruct it.
The script gives the same two means \tagC.
\end{worked}

At full size the trim is large in absolute terms. A batch of 64 fully valid patches has $64\times45{,}872=2{,}935{,}808$
hidden pixels; the trim keeps $\lfloor2{,}935{,}808\times0.95\rfloor=2{,}789{,}017$ and drops 146,791 \tagC. The 5\% is
taken over the whole batch, not per patch, so a patch with a large fire can lose more than 5\% of its pixels and a
quiet one none.

\subsection{Optimiser, schedule, data}

The optimiser is Adam at a learning rate of $10^{-4}$, lowered along half a cosine wave over 500 epochs to $10^{-6}$
(both are derived in \cref{ch:trainloop}); the configuration's
\lstinline!warmup_epochs: 5! is read by nothing in \file{app/abstract_classes/foundation_trainer.py}, which builds
\code{CosineAnnealingLR} with only \code{T_max}, \code{eta_min} and \code{last_epoch}. An epoch is a budget of 30,000
patches of $256\times256$ from the Landsat archive, $30{,}000/64=468.75$ batches, and validation draws 1,000
\tagM{configs/seg_former_exp_8.json}. By the cosine formula, epoch 495 runs at
$10^{-6}+(10^{-4}-10^{-6})\tfrac{1+\cos(495\pi/500)}{2}=1.02\times10^{-6}$ \tagC: the shipped checkpoint was taken
when the learning rate was within 2.5\% of its floor.

\subsection{Validation, and why train loss is lower}

Validation (\code{compute_validation_loss}) masks a random half of the valid tokens, then the complementary half,
keeps each pixel's prediction from the pass that hid it, and averages $|\hat T-T|$ over every valid pixel with no
trim. So the two numbers in the checkpoint measure different things \tagMr{checkpoint \texttt{train\_loss},
\texttt{avg\_val\_loss}}:
\begin{center}
train loss 0.1730\,°C (70\% hidden, trimmed)\qquad validation loss 0.2565\,°C (50\% hidden twice, untrimmed).
\end{center}
The validation number is the more honest one: a mean absolute temperature error of about a quarter of a degree on
held-out patches.

%======================================================================
\section{Inference on a scene}\label{sec:chakshu:infer}

The inferencer is \file{app/foundation_models/inferencers/segformer_mae_inferencer.py}; \cref{ch:inference} covers the
same machinery for Indradhanu. The product's settings for Chakshu come from the resolver's capability entry and
\code{InferenceConfig}'s defaults:
\begin{center}\small
\begin{tabular}{@{}ll@{}}
\toprule
setting & value\\
\midrule
patch, stride & 256, 128 (half-overlapping tiles)\\
batch & 8 patches\\
masking & checkerboard of single stage-1 tokens, two complementary passes\\
erosion of the validity mask & kernel 15: a 7-pixel margin at every invalid edge\\
skip a patch if valid share is below & 0.1\\
uncovered pixels & take the input value, so their residual is exactly 0\\
\bottomrule
\end{tabular}
\end{center}

\begin{worked}[label={we:chakshu:tiles}]{Tiles over two scene sizes}
Along an axis of length $W$ with patch $P=256$ and stride $S=128$, tiles start every $S$ pixels while a whole tile
fits, and one more is snapped back to end exactly at the edge, so the count is $\lceil(W-P)/S\rceil+1$ when
$(W-P)/S$ is not whole (\cref{eq:inference:count} derives it). For $W=1{,}000$:
\begin{align*}
(1{,}000-256)/128 &= 744/128 = 5.8125, \qquad \lceil5.8125\rceil+1 = 6+1 = 7,
\end{align*}
starts 0, 128, \dots, 640, and the snapped last one at $1{,}000-256=744$. A $1{,}000\times1{,}000$ scene needs
$7\times7=49$ tiles and $2\times49=98$ forward passes. A full Landsat scene of about $7{,}801\times7{,}681$ pixels
\tagA{} needs $60\times60=3{,}600$ tiles and 7,200 passes. The patch planner returns 49 and 3,600 \tagC.
\end{worked}

\begin{sourcenote}[: two training-inference mismatches, already known]
\file{docs/09-known-issues.md} (item 3) records both, and the source confirms them. Inference defaults to the
checkerboard while training and validation use random masks; and the inferencer's \lstinline!MIN_VALID_FRACTION = 0.1!
sits under a comment saying it ``Matches training'', where training discards below 0.4. A patch between 10\% and
40\% valid is therefore reconstructed by a model that never trained on one.
\end{sourcenote}

%======================================================================
\section{From reconstruction to score}\label{sec:chakshu:score}

The resolver offers Chakshu two scoring methods, \code{L1} (the default) and \code{MSE}. With one band both are simple
(\file{app/utils/anomaly_detection/scoring.py}, \code{compute_score}):
\begin{equation}\label{eq:chakshu:score}
s_\text{L1}(p) = \big|T(p)-\hat T(p)\big|, \qquad s_\text{MSE}(p) = \big(T(p)-\hat T(p)\big)^2,
\end{equation}
in °C and °C$^2$. There is no spectral angle (one number has no direction) and no residual Mahalanobis; there is no
segmentation head, so nothing is excluded after scoring. A \code{cloud_mask} action upstream supplies a keep mask,
and pixels outside it score 0.

\subsection{A synthetic test, with the real weights}

To see the score at work, the script builds a smooth $256\times256$ field, $T=25+3\sin(x/40)+2\cos(y/30)$\,°C, and
plants hot spots at 45\,°C in copies of it \tagA. It runs the product's inferencer on each with the shipped weights
and the product's checkerboard. These are synthetic numbers, not a measurement on Landsat scenes.
\begin{center}\small
\begin{tabular}{@{}lccc@{}}
\toprule
field & median $|T-\hat T|$, background & at the spot & largest\\
\midrule
plain & 0.019\,°C & -- & 0.092\,°C\\
1 pixel at 45\,°C & 0.019\,°C & 20.96\,°C & 20.96\,°C\\
$4\times4$ spot inside one token & 0.019\,°C & 21.05\,°C (centre) & 21.17\,°C\\
$3\times3$ spot across four tokens & 0.019\,°C & 7.59\,°C (centre) & 12.60\,°C\\
\bottomrule
\end{tabular}
\end{center}
\tagC. The background is reconstructed to two hundredths of a degree. A one-pixel or token-sized hot spot is
reconstructed as background: at the token-sized spot the model predicts 23.95\,°C where the field reads 45, and the
residual, 21.05\,°C, is the whole of the anomaly. The $3\times3$ spot is the surprise.

\subsection{Why a smaller spot can score lower}\label{sec:chakshu:straddle}

Stage-1 tokens are $4\times4$ blocks starting at multiples of 4. The $3\times3$ spot occupies rows and columns 126 to
128: rows 126 and 127 lie in token row 31 (pixels 124--127), row 128 in token row 32 (128--131), and the same for
columns. So the spot straddles a $2\times2$ group of tokens. The checkerboard colours a token by
$(\text{row}+\text{column}) \bmod 2$:
\begin{align*}
(31,31) &: 62 \bmod 2 = 0, & (31,32) &: 63 \bmod 2 = 1,\\
(32,31) &: 63 \bmod 2 = 1, & (32,32) &: 64 \bmod 2 = 0 .
\end{align*}
In the pass that hides colour 0, tokens $(31,32)$ and $(32,31)$ are visible, and each holds two of the spot's hot
pixels. The model sees part of the spot beside the hidden part and predicts warm. In the other pass the roles swap.
So every hot pixel is predicted from context that contains other pixels of the same spot: the spot helps
reconstruct itself, and the residual drops from about 21\,°C to 7.6\,°C at the centre. \Cref{fig:chakshu:straddle}
draws it.

\begin{figure}[htbp]\centering
\begin{tikzpicture}[x=5mm,y=-5mm]
  % 8x8 pixels = tokens (31,31),(31,32),(32,31),(32,32); pixel 124..131
  \foreach \tr/\tc/\col in {0/0/0,0/1/1,1/0/1,1/1/0}{
    \ifnum\col=0 \fill[maskgrey!45] (4*\tc,4*\tr) rectangle ++(4,4);
    \else \fill[spec!18] (4*\tc,4*\tr) rectangle ++(4,4); \fi}
  \fill[loss!75] (2,2) rectangle (5,5);
  \foreach \i in {0,...,8}{\draw[gray!60,line width=0.2pt] (\i,0)--(\i,8); \draw[gray!60,line width=0.2pt] (0,\i)--(8,\i);}
  \foreach \i in {0,4,8}{\draw[black,line width=0.9pt] (\i,0)--(\i,8); \draw[black,line width=0.9pt] (0,\i)--(8,\i);}
  \node[font=\scriptsize,above] at (2,0) {col 31}; \node[font=\scriptsize,above] at (6,0) {col 32};
  \node[font=\scriptsize,left] at (0,2) {row 31}; \node[font=\scriptsize,left] at (0,6) {row 32};
  \node[font=\scriptsize,anchor=west,align=left] at (8.6,1.5) {grey: hidden in pass 1\\(colour 0)};
  \node[font=\scriptsize,anchor=west,align=left] at (8.6,4) {teal: visible in pass 1\\(colour 1)};
  \node[font=\scriptsize,anchor=west,align=left] at (8.6,6.5) {amber: the $3\times3$ hot spot,\\pixels 126--128};
\end{tikzpicture}
\caption{An $8\times8$-pixel window around the $3\times3$ spot of \cref{sec:chakshu:straddle}. Thick lines are stage-1
token borders. In pass 1 the grey tokens are hidden, but the two teal tokens still show four of the spot's nine hot
pixels to the encoder, so the hidden hot pixels are predicted warm; pass 2 swaps the colours.}
\label{fig:chakshu:straddle}
\end{figure}

\begin{keypoint}
A hot object larger than one pixel can reconstruct itself through the visible tokens it reaches into; Chakshu's
residual is largest for objects that fit inside one $4\times4$ token.
\end{keypoint}

This is the analogy's large fire, and it is the same leakage Indradhanu has. It is a property of the two-pass
design, not a bug in the code, but it means a residual's size depends on where an object falls on the token grid.
Larger checkerboard cells (\code{checkerboard_cell_size}, \cref{ch:inference}) would hide $2\times2$ groups of tokens
together; the product does not use them for Chakshu, and no measurement of that option on thermal scenes exists.

%======================================================================
\section{The HotSat-1 units caveat}\label{sec:chakshu:hotsat}

HotSat-1's L2 Visual product carries uncalibrated digital numbers, ``relative radiance differences within a
scene'', around 5,000 with a spread of about 400 (the builder's and the override's docstrings,
\file{app/utils/dataset_builder/hotsat_dataset_builder.py}, \file{app/models/training/inference_config.py}). The
product scores HotSat-1 with Chakshu by default (\file{backend/allotrope/action_types/anomaly_scoring.py}), and its
plan is a \emph{per-scene override}: replace $\mu_T$ and $\sigma_T$ by the scene's own mean and spread, so the input
arrives centred with unit spread and scores become comparable within the scene, not across scenes.

\begin{worked}[label={we:chakshu:hotsat}]{A HotSat number under the baked statistics}
\begin{align*}
z &= (5{,}000-24.5756)/13.5744 = 4{,}975.4244/13.5744 = 366.53
\end{align*}
Training inputs sat within a few units of 0. The course and \file{docs/lld/pixel-stats.md} give 471.0 for this,
from a mean of 290 and spread of 10 (\cref{sec:chakshu:input}); with the real buffers it is 366.5 \tagC. Under the
override with mean 5,000 and spread 400, the same pixel gives $(5{,}000-5{,}000)/400=0$ and a pixel at 5,400 gives
$400/400=1$.
\end{worked}

That is the plan. The code defeats it twice, and the second time is new.

\begin{sloppypar}\noindent\textbf{First: the units never arrive.} The override is built only when the vendable's \code{units} begins with
\code{"DN_"}. \code{HotSatDatasetBuilder} passes \code{units="DN_14bit_relative"} to \code{VendableThermalDataset},
which declares no such field, so Pydantic discards it, and \lstinline!getattr(vendable, "units", None)! is \code{None}
(\file{docs/09-known-issues.md}, item 7; \cref{ch:product,ch:bandpipeline}).\end{sloppypar}

\begin{sloppypar}\noindent\textbf{Second: the checkpoint overwrites the override.} Suppose the units did arrive. The handler builds
\lstinline!PixelStatsOverride(mean=[5000], std=[...])!, and the inferencer constructs \code{SegFormerMAE} with those
numbers in its \code{PixelNormalize} buffers. Then \code{FoundationInferencer._load_weights} calls
\code{load_state_dict(ckpt["model_state_dict"])}, and the checkpoint's state dictionary contains
\code{normalize.mean}, \code{normalize.std}, \code{denormalize.mean} and \code{denormalize.std}: it writes 24.5756 and
13.5744 back over the override. The script shows it directly: a model built with mean 5,000.03 has mean 24.5756
after the load, and an inferencer constructed with the override holds $[24.5756]$, $[13.5744]$ \tagC.
\file{docs/lld/pixel-stats.md}'s invariant, ``The override wins outright'', holds in \code{resolve_pixel_stats} and
is undone two calls later.\end{sloppypar}

\begin{sourcenote}[: the override cannot work for any checkpoint with normalisation buffers]
Measured on CPU with the shipped weights (\file{tools/compute_chakshu.py}): a synthetic HotSat-like field
(5,000 DN, spread 193) reconstructs to a mean of 100.3 DN, a median residual of 4,899.5 DN everywhere, with and
without the override; the two runs are identical \tagC. Every pixel looks equally anomalous, which is the failure the
LLD describes as the reason the override exists. The same load-after-build order applies to every model whose
checkpoint carries \code{normalize.*} buffers: Pratibimba, Antardhana, Tirohita and Drashta (\cref{ch:convae}).
A third, smaller fault: the handler's check that the override's length matches \code{in_channels} reads
\lstinline!getattr(m.model_config, "in_channels", None)!, but \code{model_config} is a dictionary, so the attribute is
never found and the check never runs. Fixing the second fault means applying the override after the weights load, or
dropping the \code{normalize.*} keys from the state dictionary before loading; this book does not change code.
\end{sourcenote}

There is no HotSat-1 scene in the local database (\file{docs/proposals/aigc-evaluation-submission.md}, row R3), so
none of this has been seen on real data. The diagnostics would still say \lstinline!normalization_mode: "baked"!, which is
true, and nothing would warn that the input was 366 standard deviations from anything the model learned.

%======================================================================
\section{Where Chakshu sits in the product}\label{sec:chakshu:product}

\begin{center}\small
\begin{tabular}{@{}ll@{}}
\toprule
resolver field (\texttt{\_CAPABILITIES}, key \texttt{segformer\_mae}) & value\\
\midrule
\texttt{foundation\_model\_name} & \texttt{segformer\_mae} $\to$ \texttt{SegFormerMAEInferencer}\\
\texttt{scoring\_methods}, default & L1, MSE; L1\\
\texttt{default\_patch\_size}, \texttt{default\_stride} & 256, 128\\
\texttt{default\_batch\_size} & 8\\
\texttt{valid\_patch\_sizes} & 64, 128, 256 (each divisible by 32)\\
\texttt{pixel\_stats\_relpath} & \texttt{constants/thermal\_pixel\_consts.json}\\
\texttt{band\_grid}, \texttt{exclude\_classes} & common165 (not checked for thermal), none\\
\bottomrule
\end{tabular}
\end{center}
\begin{sloppypar}\noindent The scoring action's sensor defaults pick Chakshu for \code{landsat9} and \code{hotsat1}
(\file{backend/allotrope/action_types/anomaly_scoring.py}). For a thermal scene the action reads the onboarding
vendable directly (no band filter), takes \code{normalized_thermal_cube} as one band, attaches a \code{cloud_mask}
output's keep mask if there is one, and writes \file{models/Chakshu/anomaly_score.tif}, the reconstruction and their
previews. The thermal chain is \code{cloud_mask} $\to$ \code{anomaly_scoring} $\to$ \code{anomaly_detection_prep} $\to$
export, with the georeference read back from the Landsat GeoTIFF at export time (\cref{ch:material}). That chain
has not been run end to end on the local stack: no thermal scene is loaded
(\file{docs/proposals/aigc-evaluation-submission.md}).\end{sloppypar}

\begin{drift}
Read 2026-10-31 from the \file{/models/segformer_mae/} volume: current checkpoint \file{segformer_mae_v0.8.0_epoch495.pt},
epoch 495, train loss 0.1730, validation loss 0.2565 (256-pixel patches), 406,496 trainable numbers (406,500 with
buffers); normalisation $24.5756\pm13.5744$\,°C; alternative v0.1.0 (epoch 10, validation 6.59877 on 128-pixel patches),
not loadable by the current code. Synthetic-field residuals in \cref{sec:chakshu:score} were computed with these
weights \tagC.
\end{drift}

%======================================================================
\begin{confusions}
\confusion{Chakshu reads Kelvin.}{It reads Celsius. The Landsat builder converts to Celsius and the buffers are
24.58 and 13.57. ``Kelvin'' and ``290'' in comments and in the pixel-stats LLD are wrong.}
\confusion{Chakshu is SegFormer-B0 with one input channel.}{The classes are, the sizes are not: widths 16 to 96, one
block per stage, expansion 2, decoder width 96. Only the unused v0.1.0 run is B0.}
\confusion{406,500 parameters.}{406,496 parameters and 4 buffers. The 406,500 is the state dictionary's size.}
\confusion{The small model beat the large one 25 times over.}{The large run had 10 epochs against 495, and their
losses were on different patch sizes. It shows the large run was unfinished.}
\confusion{Train loss 0.173 below validation 0.257 means overfitting.}{The two measure different things: training
trims the worst 5\% and hides 70\%; validation hides half twice and trims nothing.}
\confusion{The per-scene override makes Chakshu safe for HotSat-1.}{It is never built (the units field is dropped),
and if it were built the checkpoint's buffers would overwrite it when the weights load.}
\confusion{A bigger hot spot always scores higher.}{Not with two-pass masking: a spot that straddles token borders
shows part of itself to the encoder in both passes and reconstructs itself partly; a spot inside one token is
hidden whole.}
\confusion{With one band, L1 and MSE rank pixels differently.}{They cannot: the square is increasing for non-negative
numbers, so $|r_1|<|r_2|$ exactly when $r_1^2<r_2^2$, and every threshold on one is a threshold on the other. They
differ only in how the map looks.}
\end{confusions}

\begin{checkyourself}
\item\label{ex:chakshu:1} Use \cref{eq:chakshu:block} to count a Chakshu stage-2 block ($C=32$, $R=4$, $E=2$), and the
same block with $E=4$. How much does doubling the expansion add, and why is it exactly $2EC^2$-shaped?
\item\label{ex:chakshu:2} At a patch size of 128 instead of 256, give the token grid, the number of tokens and the
number of attention keys at each stage, and the number of tokens a training pass hides at $\rho=0.7$. Does the
parameter count change?
\item\label{ex:chakshu:3} A Landsat pixel has DN $=48{,}000$. Convert it to °C with \cref{eq:chakshu:dn} and normalise
it with \cref{eq:chakshu:z}.
\item\label{ex:chakshu:4} Ten hidden pixels have errors 0.1, 0.3, 0.2, 0.1, 0.4, 0.2, 0.1, 0.3, 0.2 and 5.0\,°C. Compute
the plain mean and the trimmed mean at $\tau=0.05$. Then at $\tau=0.1$. What does the floor do at $\tau=0.05$ with
only ten numbers?
\item\label{ex:chakshu:5} A scene is $1{,}000\times600$ pixels. How many Chakshu tiles does the product run, and how
many forward passes?
\item\label{ex:chakshu:6} A HotSat-1 pixel reads 4,600 DN. Compute its $z$ under the baked statistics, and under a
per-scene override of mean 5,000 and spread 400. Which $z$ does the product actually use today, and why?
\item\label{ex:chakshu:7} A $2\times2$ hot spot covers pixels 127--128 in both rows and columns. Which stage-1 tokens does
it touch, what are their checkerboard colours, and would you expect its residual to be closer to the one-token spot's
or the $3\times3$ spot's?
\item\label{ex:chakshu:8} The resolver hard-codes heads, ratios, blocks and expansion for \code{segformer_mae}. Name
the one manifest field that changes the model's width and the one that changes the decoder, and say what would happen
if a two-blocks-per-stage thermal checkpoint were published with the same manifest format.
\end{checkyourself}

\nextchapter{ch:ceiling}
